\section{Discussion and conclusion}
\label{sec:discussion}

Previous dissipative QRC studies largely fix the environmental process and tune
only its strength. We show that this one-rate picture is incomplete. In the
finite spin reservoirs studied here, whether qubit transitions couple to
separate or shared channels, and how coupling is distributed among them,
changes which parts of the input history the measured output can recover. A
reservoir's connection to the environment is therefore part of its
computational design.

The clearest recurring comparison is between independent and shared
relaxation. With independent relaxation, each qubit couples to its own decay
channel. With shared relaxation, the environment couples directly to one
collective combination of transitions while the Hamiltonian can mix other
combinations into it. Shared relaxation makes more recent input history
recoverable and improves a nonlinear benchmark that depends on several earlier
inputs.

The cleanest intervention changes only the collective combination addressed by
an otherwise matched shared channel. Memory changes even though the channels
have the same number of coupled directions, the same nonzero coupling
strengths, the same total and per-qubit coupling weights, and identical
processors, inputs, and readouts. The same direction sensitivity appears at
both tested sizes. These comparisons identify what coarse summaries of
dissipation miss: the collective transition directly coupled to the
environment. Dynamical diagnostics and a gradual path between independent and
shared relaxation support this channel-selectivity picture, but do not
establish a unique microscopic mechanism.

The local-versus-shared ordering persists after longer washout, across the
tested sizes and Hamiltonians, with a different input encoding, and under
separate controls for average dissipative activity, a representative relaxation
rate, and operating-point selection. These checks make incomplete washout, one
processor or encoding, and the matched scalar differences insufficient
explanations. They establish a recurring contrast, not a universal advantage
over every environmental process.

Across memory, nonlinear processing, parity, and chaotic forecasting, different
dissipative designs favor different tasks. Redistributing a fixed coupling
budget within one design also changes memory. No process is best for every task.
Coupling organization is therefore a task-dependent design variable, not a
universal optimization rule.

Memory here means previous-input information that is linearly recoverable from
the specified Pauli measurements after washout, not an intrinsic quantum-memory
lifetime. The full task map concerns numerical five-qubit driven-spin
reservoirs. Recurrence from four to eight qubits does not establish asymptotic
scaling, and the four-qubit boundary is excluded from claims that the initial
state has been fully forgotten. Holding one common measure of coupling weight
fixed provides a controlled comparison, not dynamical, physical, or hardware
equivalence. The finite-measurement analysis includes projection noise and the
chosen estimators, but not drift, gate and readout errors, latency, or full
experimental cost. The analytical result on forgetting applies to driven
dephasing, not the shared-channel effect, and a generic common bath need not
realize the modeled collective process.

The result brings a familiar principle of dissipation engineering into temporal
quantum computation: designing the interaction between the system and the
environment~\cite{Poyatos1996,Verstraete2009,Pastawski2011,Harrington2022}. Shared
resonators, waveguides, and structured couplers suggest possible routes for
testing this coupling organization
~\cite{Cattaneo2019,Brehm2021,Sheremet2023,CattaneoPRX2023}. The next challenge
is to predict useful coupling patterns jointly from the Hamiltonian, input
encoding, environment, and measurement, then compare realizable couplers under
meaningful physical budgets in larger, noisy reservoirs. How coupling
organization interacts with non-Markovian memory is a separate open question.

Together, these results move dissipation in quantum reservoir computing from a
parameter to be tuned to a connection to be designed. The question is not only
how quickly information should decay, but which transitions should share a
route to the environment so that task-relevant history remains readable. For
temporal quantum processing, the way a reservoir forgets is part of how it
computes.
